\section{Discussion and conclusion}
\label{sec:conclusion}

\subsection{Limitations}

Three limitations are worth stating. First, \cref{prop:conflict} assumes
uniformly placed edits; real revisions cluster around the parts of a paper an
author is working on, which by \cref{cor:clustered} raises the conflict rate
and argues for fewer, longer tasks over many short ones. Second,
\cref{alg:settle} cannot settle formulas, since their printed form shares no
letters with their source; for those, SyncTeX's lines are kept and the excerpt
shows the whole display. Third, an agent that cannot ask for permission, such
as a headless command line, must be restricted in advance; Papilot runs such
agents so that they edit only the task's worktree, which also means they cannot
compile or test their own changes.

\subsection{Conclusion}

Papilot treats agent editing as a sequence of reviewable rounds. Isolation
makes proposals safe, typeset review makes them legible, and explicit
acceptance keeps the author in charge of the paper. The models in this paper
make those properties checkable: compatible edits commute and rebase
(\cref{prop:commute,lem:tp1}), tasks need few rounds
(\cref{prop:rounds,cor:rounds}), undo never recompiles
(\cref{prop:undo-build}), per-word text bounds the selection error
(\cref{thm:per-word}), and partial commits keep unrelated staged work intact
(\cref{prop:partial}). They also give the reader something worth reviewing.
